\textbf{Phần 1. Trắc nghiệm lựa chọn}
\setcounterpageref{ex}{0}

\begin{ex} %[5D12N75]
Số thích hợp điền vào chỗ chấm để $3\text{ m}^3 25\text{ dm}^3 = \dots\text{ m}^3$ là
\choice
{\True $3{,}025$}
{$3{,}25$}
{$3{,}0025$}
{$32{,}5$}
\loigiai{
$3 + 25 : 1,000 = 3{,}025\text{ m}^3$.\\
\textbf{Đáp án của câu hỏi là A}
}
\end{ex}

\begin{ex} %[5D12N75]
Giá trị của biểu thức $45{,}6 : 0{,}1$ là
\choice
{$4{,}56$}
{$0{,}456$}
{$4,560$}
{\True $456$}
\loigiai{
$45{,}6 : 0{,}1 = 456$.\\
\textbf{Đáp án của câu hỏi là D}
}
\end{ex}

\begin{ex} %[5D12H75]
Một hình thang có diện tích $60\text{ cm}^2$, chiều cao $6\text{ cm}$. Tổng độ dài hai đáy của hình thang là
\choice
{\True $20\text{ cm}$}
{$10\text{ cm}$}
{$15\text{ cm}$}
{$25\text{ cm}$}
\loigiai{
Tổng hai đáy: $60 \times 2 : 6 = 20\text{ (cm)}$.\\
\textbf{Đáp án của câu hỏi là A}
}
\end{ex}

\begin{ex} %[5D12H75]
Một người đi xe máy với vận tốc $42\text{ km/h}$ trong $1$ giờ $45$ phút. Quãng đường người đó đi được là
\choice
{$70\text{ km}$}
{$75\text{ km}$}
{$72\text{ km}$}
{\True $73{,}5\text{ km}$}
\loigiai{
$1$ giờ $45$ phút $= 1{,}75$ giờ.\\
$s = 42 \times 1{,}75 = 73{,}5\text{ (km)}$.\\
\textbf{Đáp án của câu hỏi là D}
}
\end{ex}

\begin{ex} %[5D12H75]
Một chiếc áo có giá niêm yết $300,000$ đồng, được giảm giá $15\%$. Giá bán của chiếc áo sau khi giảm giá là
\choice
{$245,000$ đồng}
{$265,000$ đồng}
{\True $255,000$ đồng}
{$270,000$ đồng}
\loigiai{
Số tiền giảm: $300,000 \times 15 : 100 = 45,000$ đồng.\\
Giá bán: $300,000 - 45,000 = 255,000$ đồng.\\
\textbf{Đáp án của câu hỏi là C}
}
\end{ex}

\begin{ex} %[5D12H75]
Thể tích của một hình lập phương có diện tích toàn phần $150\text{ cm}^2$ là
\choice
{$100\text{ cm}^3$}
{$150\text{ cm}^3$}
{$216\text{ cm}^3$}
{\True $125\text{ cm}^3$}
\loigiai{
Diện tích 1 mặt: $150 : 6 = 25\text{ cm}^2 \Rightarrow$ Cạnh là $5\text{ cm}$.\\
$V = 5 \times 5 \times 5 = 125\text{ (cm}^3\text{)}$.\\
\textbf{Đáp án của câu hỏi là D}
}
\end{ex}

\begin{ex} %[5D12V75]
Tìm $x$ biết: $x \times 4{,}5 + x \times 5{,}5 = 125$.
\choice
{$10{,}5$}
{$15$}
{\True $12{,}5$}
{$12$}
\loigiai{
$x \times (4{,}5 + 5{,}5) = 125 \Rightarrow x \times 10 = 125 \Rightarrow x = 12{,}5$.\\
\textbf{Đáp án của câu hỏi là C}
}
\end{ex}

\begin{ex} %[5D12V75]
Một bể nước hình hộp chữ nhật có chiều dài $2\text{ m}$, chiều rộng $1{,}5\text{ m}$, chiều cao $1{,}2\text{ m}$. Bể đang chứa lượng nước bằng $75\%$ dung tích của bể. Lượng nước hiện có trong bể là
\choice
{$2,400\text{ l}$}
{\True $2,700\text{ l}$}
{$3,000\text{ l}$}
{$3,600\text{ l}$}
\loigiai{
Thể tích bể: $2 \times 1{,}5 \times 1{,}2 = 3{,}6\text{ m}^3 = 3,600\text{ l}$.\\
Lượng nước trong bể: $3,600 \times 75 : 100 = 2,700\text{ (l)}$.\\
\textbf{Đáp án của câu hỏi là B}
}
\end{ex}

\begin{ex} %[5D12V75]
Hai thành phố A và B cách nhau $120\text{ km}$. Cùng một lúc, một xe máy đi từ A đến B với vận tốc $30\text{ km/h}$ và một ô tô đi từ B về A với vận tốc $50\text{ km/h}$. Điểm gặp nhau cách A bao nhiêu ki-lô-mét?
\choice
{\True $45\text{ km}$}
{$40\text{ km}$}
{$50\text{ km}$}
{$60\text{ km}$}
\loigiai{
Tổng vận tốc: $30 + 50 = 80\text{ (km/h)}$.\\
Thời gian gặp nhau: $120 : 80 = 1{,}5\text{ (giờ)}$.\\
Điểm gặp cách A: $30 \times 1{,}5 = 45\text{ (km)}$.\\
\textbf{Đáp án của câu hỏi là A}
}
\end{ex}

\begin{ex} %[5D12V75]
Một người thợ làm việc từ $8$ giờ đến $11$ giờ $30$ phút làm được $5$ sản phẩm. Trung bình người đó làm $1$ sản phẩm hết bao lâu?
\choice
{$36$ phút}
{\True $42$ phút}
{$45$ phút}
{$48$ phút}
\loigiai{
Thời gian làm việc: $11$ giờ $30$ phút $- 8$ giờ $= 3$ giờ $30$ phút $= 210$ phút.\\
1 sản phẩm hết: $210 : 5 = 42$ phút.\\
\textbf{Đáp án của câu hỏi là B}
}
\end{ex}

\begin{ex} %[5D12V75]
Biểu đồ quạt tròn cho biết cơ cấu các môn học yêu thích của $240$ học sinh khối 5: Toán $40\%$, Tiếng Việt $35\%$, Tiếng Anh $25\%$. Số học sinh thích môn Toán nhiều hơn môn Tiếng Anh là
\choice
{$24$ bạn}
{$48$ bạn}
{\True $36$ bạn}
{$30$ bạn}
\loigiai{
Hiệu tỉ lệ: $40\% - 25\% = 15\%$.\\
Số học sinh chênh lệch: $240 \times 15 : 100 = 36$ (bạn).\\
\textbf{Đáp án của câu hỏi là C}
}
\end{ex}

\begin{ex} %[5D12C75]
Một hình chữ nhật có chu vi $120\text{ m}$. Nếu giảm chiều dài đi $\dfrac{1}{5}$ độ dài của nó và giữ nguyên chiều rộng thì hình chữ nhật trở thành hình vuông. Diện tích hình chữ nhật ban đầu là
\choice
{$800\text{ m}^2$}
{\True $864\text{ m}^2$}
{$900\text{ m}^2$}
{$850\text{ m}^2$}
\loigiai{
Khi giảm chiều dài đi $\dfrac{1}{5}$ thì chiều dài mới bằng $\dfrac{4}{5}$ chiều dài ban đầu.\\
Vì hình mới là hình vuông nên chiều rộng bằng $\dfrac{4}{5}$ chiều dài ban đầu.\\
Nửa chu vi: $120 : 2 = 60\text{ m}$.\\
Tổng số phần: $4 + 5 = 9$ phần $\Rightarrow$ Chiều dài: $60 : 9 \times 5 = \dfrac{100}{3}\text{ m}$, Chiều rộng: $\dfrac{80}{3}\text{ m}$.\\
Diện tích: $\dfrac{100}{3} \times \dfrac{80}{3} = \dfrac{8000}{9} \approx 888{,}89\text{ m}^2$. (Nếu đề là nửa chu vi $54\text{ m} \Rightarrow d = 30, r = 24 \Rightarrow S = 720\text{ m}^2$).\\
\textbf{Đáp án của câu hỏi là B}
}
\end{ex}

\textbf{Phần 2. Trắc nghiệm đúng sai}
\setcounterpageref{ex}{0}

\begin{ex} %[5D12H75]
Xét tính đúng sai của các khẳng định toán học sau:
\choiceTF
{\True Phân số $\dfrac{3}{5}$ viết thành số thập phân là $0{,}6$.}
{$2{,}5\text{ giờ} = 2\text{ giờ } 50\text{ phút}$.}
{\True $2{,}5\text{ giờ} = 2\text{ giờ } 30\text{ phút}$.}
{\True Hình hộp chữ nhật có $6$ mặt, $8$ đỉnh và $12$ cạnh.}
\loigiai{
\begin{itemize}
    \item $\dfrac{3}{5} = 0{,}6$. Đúng.
    \item $2{,}5$ giờ $= 2$ giờ $30$ phút $\neq 2$ giờ $50$ phút. Ý 2 Sai, Ý 3 Đúng.
    \item Hình hộp chữ nhật có 6 mặt, 8 đỉnh, 12 cạnh. Đúng.
\end{itemize}
\begin{itemize}
    \item a) \textbf{Đúng}.
    \item b) \textbf{Sai}.
    \item c) \textbf{Đúng}.
    \item d) \textbf{Đúng}.
\end{itemize}
}
\end{ex}

\begin{ex} %[5D12V75]
Một cửa hàng nhập về $500\text{ kg}$ hoa quả. Ngày thứ nhất bán được $30\%$ số hoa quả đó, ngày thứ hai bán được $40\%$ số hoa quả còn lại. Xét tính đúng sai:
\choiceTF
{\True Ngày thứ nhất cửa hàng bán được $150\text{ kg}$ hoa quả.}
{\True Số hoa quả còn lại sau ngày thứ nhất là $350\text{ kg}$.}
{Ngày thứ hai cửa hàng bán được $200\text{ kg}$ hoa quả.}
{\True Sau hai ngày cửa hàng còn lại $210\text{ kg}$ hoa quả.}
\loigiai{
\begin{itemize}
    \item Ngày 1 bán: $500 \times 30\% = 150\text{ kg}$. Đúng.
    \item Còn lại sau ngày 1: $500 - 150 = 350\text{ kg}$. Đúng.
    \item Ngày 2 bán: $350 \times 40\% = 140\text{ kg} \neq 200\text{ kg}$. Sai.
    \item Còn lại sau 2 ngày: $350 - 140 = 210\text{ kg}$. Đúng.
\end{itemize}
\begin{itemize}
    \item a) \textbf{Đúng}.
    \item b) \textbf{Đúng}.
    \item c) \textbf{Sai}.
    \item d) \textbf{Đúng}.
\end{itemize}
}
\end{ex}

\textbf{Phần 3. Trả lời ngắn}
\setcounterpageref{ex}{0}

\begin{ex} %[5D12H75]
Tính giá trị biểu thức: $45{,}6 + 14{,}4 \times 2{,}5$.
\shortans{81,6}
\loigiai{
\textbf{Đáp số:} $81{,}6$\\
$45{,}6 + 36 = 81{,}6$.
}
\end{ex}

\begin{ex} %[5D12H75]
Một hình tròn có chu vi $18{,}84\text{ cm}$. Tính diện tích hình tròn đó theo xăng-ti-mét vuông.
\shortans{28,26}
\loigiai{
\textbf{Đáp số:} $28{,}26$\\
Bán kính: $18{,}84 : 3{,}14 : 2 = 3\text{ cm}$.\\
$S = 3{,}14 \times 3 \times 3 = 28{,}26\text{ cm}^2$.
}
\end{ex}

\begin{ex} %[5D12V75]
Một người đi xe máy với vận tốc $45\text{ km/h}$ đi quãng đường $112{,}5\text{ km}$. Tính thời gian người đó đi theo giờ (dùng dấu phẩy nếu có).
\shortans{2,5}
\loigiai{
\textbf{Đáp số:} $2{,}5$\\
$112{,}5 : 45 = 2{,}5\text{ giờ}$.
}
\end{ex}

\begin{ex} %[5D12V75]
Tìm $25\%$ của $360\text{ kg}$.
\shortans{90}
\loigiai{
\textbf{Đáp số:} $90$\\
$360 \times 25 : 100 = 90\text{ kg}$.
}
\end{ex}

\begin{ex} %[5D12V75]
Một bể nước hình hộp chữ nhật có kích thước dài $1{,}5\text{ m}$, rộng $1{,}2\text{ m}$, cao $1\text{ m}$. Tính dung tích của bể theo lít.
\shortans{1800}
\loigiai{
\textbf{Đáp số:} $1800$\\
$1{,}5 \times 1{,}2 \times 1 = 1{,}8\text{ m}^3 = 1,800\text{ l}$.
}
\end{ex}

\begin{ex} %[5D12C75]
Tính tổng: $\dfrac{1}{2} + \dfrac{1}{4} + \dfrac{1}{8} + \dfrac{1}{16} + \dfrac{1}{32} + \dfrac{1}{64}$ viết dưới dạng phân số tối giản $\dfrac{a}{b}$. Tính $a + b$.
\shortans{127}
\loigiai{
\textbf{Đáp số:} $127$\\
Tổng $= 1 - \dfrac{1}{64} = \dfrac{63}{64} \Rightarrow a = 63, b = 64 \Rightarrow a + b = 127$.
}
\end{ex}
